\documentclass[12pt]{article}
\usepackage[utf8]{inputenc}
\usepackage[T1]{fontenc}
\usepackage{hyperref}
\title{Minkin Language Resources for K-12 Curriculum Development}
\author{Compiled by Claude Sonnet 5.5 (claude-sonnet-5-5) for Andr\'as Kornai}
\date{}
\begin{document}
\maketitle

\section{Introduction}
Minkin (also Mingginda, Minginnda, Mingin; ISO 639-3 xxm; Glottocode mink1237) is an extinct Aboriginal language of the Burketown area, Queensland, on the southern side of the Gulf of Carpentaria. It has been extinct since the late 19th century. Documentation is thin: historical wordlists from the 1880s and Nicholas Evans's 1990 article ``The Minkin language of Burketown''. Evans and Bowern (2011) treat it as a Tangkic language. Others have proposed Yiwaidjan links or an isolate. Its neighbours spoke closely related languages, including Yukulta and Ganggalida.

\section{Primary Educational Resources}
No primers, textbooks or school programmes for Minkin were found in the search results.

\section{Linguistic and Reference Materials}
\begin{itemize}
\item Wikipedia, Minkin language: \url{https://en.wikipedia.org/wiki/Minkin_language}
\item Wikipedia, Mingginda language: \url{https://en.wikipedia.org/wiki/Mingginda_language}
\item Wikipedia, Tangkic languages (mirror copy): \url{https://static.hlt.bme.hu/semantics/external/pages/jokuc/en.wikipedia.org/wiki/Tangkic_languages.html}
\item Wiktionary entry, Minkin: \url{https://en.wiktionary.com/wiki/Minkin}
\end{itemize}

\section{Religious and Community Resources}
\begin{itemize}
\item Wikipedia, Mingin people: \url{https://en.wikipedia.org/wiki/Mingin_people}
\item State Library of Queensland, International Year of Indigenous Languages word-week blog post: \url{https://www.slq.qld.gov.au/blog/2019-international-year-indigenous-languages-word-week-week-forty-nine}
\end{itemize}
No scripture translations were found.

\section{Digital and Audio Resources}
No audio, corpus or software resources were found in the search results.

\section{Curriculum Development Context}
There are no speakers and only a few historical words survive, so a curriculum taught in Minkin is not possible. A school could offer a heritage-awareness unit on the Mingin people and the Burketown region, and on neighbouring languages that are better documented. Teachers should work with Aboriginal community members and local language centres. Evans (1990) is the main scholarly source and should be consulted directly.
\end{document}